\begin{solapartat}
El període d'oscil·lació d'una molla ve donat per l'expressió:
\[ T = 2\pi\sqrt{\frac{m}{k}}\ \text{\punts{0,2}} \Rightarrow m = \frac{k}{4\pi^2}\,T^2 \Rightarrow k = \frac{4\pi^2\,m}{T^2}\ \text{\punts{0,2}} \]

Per tant si llegim un valor de $m$ i el seu valor corresponent de $T^2$ sobre la recta, obtindrem el valor de $k$:
\[ k = \frac{4\pi^2\,1,9\ \mathrm{kg}}{0,5\ \mathrm{s^2}} = 150\ \mathrm{N/m}\ \text{\punts{0,2}} \]

L'energia total és:
\[ E_T = \frac{1}{2}\,k\,A^2\ \text{\punts{0,2}} = 0,75\ \mathrm{J} \]

L'energia cinètica màxima l'obtindrem quan la seva energia potencial sigui zero i en aquest cas serà igual a l'energia total: $\Rightarrow$
\[ E_{c_\textit{màxima}} = 0,75\ \mathrm{J}\ \text{\punts{0,2}} \]
\end{solapartat}

\begin{solapartat}
\begin{gather*}
\omega = \sqrt{\frac{k}{m}} = \sqrt{\frac{150}{1,5}} = 10\ \mathrm{rad/s}\ \text{\punts{0,2}}\\
a_\textit{maàxima} = A\,\omega^2 = 20\ \mathrm{m/s^2}\ \text{\punts{0,3}} % DUBTE: «maàxima» (sic) a la pauta original
\end{gather*}

En aquest cas l'energia total de la oscil·lació és:
\[ E_T = \frac{1}{2}\,k\,A^2 = 3\ \mathrm{J}\ \text{\punts{0,2}} \]

Per parar la oscil·lació la força de fregament farà un treball igual a l'energia total de la oscil·lació:
\[ W_\textit{fregament} = 3\ \mathrm{J}\ \text{\punts{0,3}} \]
\end{solapartat}
